% section_5_conclusion.tex — v6 from-scratch rewrite
% Spec: docs/Draft/THESIS_SKELETON.md v6 §5 outline
% Phase 8 — synthesis of §3 (cash) + §4.1 (drivers) + §4.2 (spread + SEC) + §4.3 (endogeneity)
% Numbers cross-referenced to body sections; no fresh empirical claims

\section{Discussion and Conclusion}
\label{sec:conclusion}

\subsection{Summary of Findings}
\label{sec:conclusion:summary}

This paper extends the three-component CEO speech-uncertainty decomposition of \citeA{dzielinski2021}---persistent CEO style \texttt{ClarityCEO}, call-level Q\&A residual \texttt{UncResCEO}, and presentation-segment uncertainty share \texttt{UncPreCEO}---from the market-reaction, analyst-behavior, firm-performance, and CEO-compensation outcome set on which \citeauthor{dzielinski2021} validate it to a class of outcomes their paper does not test: firm-level financing-policy decisions and outsider regulatory reaction. On a panel of $112{,}968$ earnings-call transcripts covering $2{,}429$ U.S.\ firms over fiscal years 2002--2018, four empirical patterns emerge.

First, the Q\&A residual \texttt{UncResCEO} loads positively on cash holdings in 12 of 12 fixed-effects specifications at $p<0.10$ one-tailed, with point estimates between $+0.0016$ and $+0.0046$ across the Full and QtrExp measurement methods. Persistent CEO style \texttt{ClarityCEO} loads negatively in 9 of 12 specifications (5 of 12 at $p<0.05$ in Full; 9 of 12 at $p<0.05$ and 6 of 12 at $p<0.01$ in QtrExp). \texttt{UncPreCEO} is null on the cash margin once the Q\&A components are jointly included. Second, the cash response amplifies among credit-constrained firms: the \texttt{UncResCEO}~$\times$~Unrated interaction is positive at $p<0.05$ in 2 of 8 cells on the lead-DV specifications, and the \texttt{UncResCEO}~$\times$~HighCFvol interaction reproduces the asymmetry through the cash-flow-volatility moderator at one tier weaker. Third, all four uncertainty drivers (firm political risk, U.S.\ economic-policy uncertainty, global economic-policy uncertainty, and product-market competition) load positively on the speech-uncertainty constructs, establishing construct validity for the decomposition before the firm-side outcome tests. Fourth, the post-call quoted bid-ask spread and SEC comment-letter receipt both load on \texttt{UncPreCEO} (4 of 12 spread cells and 4 of 6 SEC cells significant at $p<0.10$); the Q\&A components are null on both outsider-reaction outcomes.

\subsection{Channel-Asymmetry Synthesis}
\label{sec:conclusion:asymmetry}

The four empirical patterns share a single structural interpretation. The firm's own financing-policy response (Section~\ref{sec:main}) loads on the Q\&A components of CEO speech; the outsider-reaction outcomes (Section~\ref{sec:additional:reaction}) load on the presentation segment. The asymmetry sits cleanly on \citeauthor{dzielinski2021}'s own theoretical structure for the two segments of the call: the presentation segment is the one \citeauthor{dzielinski2021} characterize as ``scripted and vetted beforehand\dots\ likely to reflect the firm's culture'', and the Q\&A segment is the one they characterize as ``inevitably somewhat improvised and hence more likely to reveal each CEO's style.'' Two segments of the same call carry signals to actors with different information sets and different decision horizons. Insiders making the firm's financing decisions track Q\&A signals---both the persistent CEO trait and the call-level state residual; outsiders running market-microstructure and regulatory-review processes track the presentation share. We document the asymmetry within \citeauthor{dzielinski2021}'s theoretical scaffolding rather than impose it.

\subsection{Mechanism Disclosure}
\label{sec:conclusion:mechanism}

The Section~\ref{sec:main} firm-side cash response and the Section~\ref{sec:additional:reaction} outsider-reaction responses operate through theoretically distinct mechanisms; we do not claim a causal bridge between them. The firm-side response is anchored in the precautionary motive of \citeA{opler1999} and \citeA{bates2009}: uncertainty raises optimal cash, with the call-level Q\&A residual a real-time within-call uncertainty signal. The outsider-reaction response is anchored in information-asymmetry and regulatory-review processes (\citeA{bushee2018} for market microstructure; \citeA{lerman2026} for SEC scrutiny), where the scripted prepared-remarks share is the natural object outsiders parse against the firm's mandatory filings. The two mechanisms can hold simultaneously without either causing the other: the same call carries Q\&A signals to insiders and presentation signals to outsiders, and both populations respond in the directions their respective theoretical models predict.

Section~\ref{sec:additional:endo} addresses three identification threats to the firm-side cash response with three orthogonal designs: a CEO sudden-death difference-in-differences against the reverse-causality threat, a first-difference replication of \citeauthor{dzielinski2021}'s Section~6 against time-invariant heterogeneity, and a \citeA{lewbel2012} heteroskedasticity-based instrumental-variable estimator against time-varying omitted confounders. The three designs target different identification threats and different speech components by construction: the Phase~E and DWZ first-difference designs identify the persistent-style component; the Lewbel design identifies the call-level residual. The asymmetry reflects the underlying statistical structure of the two speech components, not a coverage gap.

\subsection{Limitations}
\label{sec:conclusion:limitations}

Three limitations qualify the empirical evidence. First, the Phase~E sudden-death sample size of eight viable events is below conventional power thresholds; the four DiD specifications report negative point estimates in the predicted direction but none statistically significant at $p<0.10$ one-tailed, and the originally proposed heterogeneity test on pre-event speech-uncertainty composition is not feasible. Second, the \citeA{lewbel2012} heteroskedasticity-based estimator delivers a 2SLS point estimate approximately five times the OLS estimate but with wider standard errors; the Wu-Hausman test does not reject OLS-consistency, and we treat this as a failure to reject rather than evidence in favor of OLS-consistency. Third, the cash dependent variable is the linear cash-to-assets ratio of \citeauthor{bates2009}; the log-of-cash-to-net-assets specification of \citeauthor{opler1999} is flagged as a planned robustness extension but not estimated in the present draft.

\subsection{Future Research}
\label{sec:conclusion:future}

Four directions extend this paper's results. First, the Phase~E sudden-death DiD heterogeneity test by pre-event \texttt{UncResCEO} composition is naturally extended to a larger event sample. Second, the \citeauthor{opler1999} log-of-cash-to-net-assets robustness is a one-runner extension reported in the body's robustness section. Third, the firm-side response to CEO speech uncertainty is documented at the cash margin; extension to other financing-policy margins (leverage, payout, capital structure) under the same three-component decomposition is a natural continuation. Fourth, the channel asymmetry between insider and outsider reactions admits sharpening through experimental or quasi-experimental settings that vary the segment of the call exogenously; we leave this as a methodological direction for the disclosure-quality literature.
